% BHU PYQ -- Auto-generated & error-corrected
\documentclass[11pt,a4paper]{article}

\usepackage[a4paper,top=2.5cm,bottom=2.5cm,left=2.8cm,right=2.8cm,headheight=14pt]{geometry}
\usepackage{amsmath,amssymb}
\usepackage[T1]{fontenc}
\usepackage[utf8]{inputenc}
\usepackage{lmodern}
\usepackage{microtype}
\usepackage{fancyhdr}
\usepackage{enumitem}
\usepackage{hyperref}
\usepackage{lastpage}
\usepackage{parskip}

\hypersetup{colorlinks=true,urlcolor=black,linkcolor=black,
  pdftitle={STB-101: Statistical Methods and Probability -- BHU PYQ},pdfauthor={Banaras Hindu University}}

\pagestyle{fancy}\fancyhf{}
\renewcommand{\headrulewidth}{0.4pt}\renewcommand{\footrulewidth}{0.4pt}
\fancyhead[L]{\small\textit{Banaras Hindu University}}
\fancyhead[R]{\small\textit{STB-101: Statistical Methods and Probability}}
\fancyfoot[C]{\small Page \thepage\ of \pageref{LastPage}}

\newlist{parts}{enumerate}{2}
\setlist[parts,1]{label=(\alph*),leftmargin=*,itemsep=5pt,topsep=3pt}
\setlist[parts,2]{label=(\roman*),leftmargin=*,itemsep=3pt,topsep=2pt}
\newcommand{\pts}[1]{\hfill\textit{[\#1]}}

\begin{document}

\begin{center}
  {\large\bfseries BANARAS HINDU UNIVERSITY}\\[2pt]
  {\normalsize B.Sc. (Hons.) Semester I Examination, 2023-24}\\[6pt]
  \rule{0.8\linewidth}{0.5pt}\\[6pt]
  {\large\bfseries Paper No. STB-101: Statistical Methods and Probability}\\[4pt]
  {\normalsize Time: 3 Hours\hfill Maximum Marks: 70}
\end{center}
\rule{\linewidth}{0.4pt}
\smallskip
\noindent\textit{Instructions: Answer \textbf{five} questions including Question~1 which is compulsory. Symbols have their usual meanings.}
\bigskip

\medskip\noindent\textbf{Question 1.}
\begin{parts}
  \item What do you mean by the measure of central tendency? Define arithmetic mean and also write down its merits and demerits. \pts{7}
  \item Derive the formula for the arithmetic mean of a composite series having two component series. Next, the average salary of male employees in a firm was Rs 5,200 and that of females in the same firm was Rs 4,200. Find the percentage of male and female employees. \pts{7}
\end{parts}

\medskip\rule{\linewidth}{0.2pt}

\medskip\noindent\textbf{Question 2.}
\begin{parts}
  \item Define various partition values. Also, explain the procedure of how these can be located graphically. \pts{7}
  \item Define standard deviation and root mean square deviation. Also, derive a relation between standard deviation and mean deviation of a series. \pts{7}
\end{parts}

\medskip\rule{\linewidth}{0.2pt}

\medskip\noindent\textbf{Question 3.}
\begin{parts}
  \item Define $r$-th moment about origin and $r$-th moment about mean. Also, derive the relation between moments about mean in terms of moments about origin. \pts{7}
  \item For a distribution, the mean is 10, variance is 16, $\gamma_1 = +1$, $\beta_2 = 4$. Obtain the first four moments about the origin. \pts{7}
\end{parts}

\medskip\rule{\linewidth}{0.2pt}

\medskip\noindent\textbf{Question 4.}
\begin{parts}
  \item Give the mathematical definition of probability. Next, let an urn contain 6 white, 4 red, and 9 black balls. If 3 balls are drawn at random, find the probability that:
  \begin{parts}
    \item two of the drawn balls are white;
    \item one is of each colour;
    \item none is red;
    \item at least one is white.
  \end{parts} \pts{7}
  \item Define conditional probability and independence of events. Also, state and prove the multiplication theorem of probability for any two independent events. \pts{7}
\end{parts}

\medskip\rule{\linewidth}{0.2pt}

\medskip\noindent\textbf{Question 5.}
\begin{parts}
  \item A and B throw alternately with a pair of balanced dice. A wins if he throws a sum of six before B throws a sum of seven; while B wins if he throws a sum of seven before A throws a sum of six. If A begins the game, then find the probability of his winning the game. \pts{7}
  \item If $A_1, A_2, \dots, A_n$ be any $n$ mutually exclusive events with $P(A_i) > 0$ for all $i = 1, 2, \dots, n$, and let $A$ be any event such that $A \subset \bigcup_{i=1}^n A_i$ with $P(A) > 0$, then find the value of $P(A_i | A)$ for all $i = 1, 2, \dots, n$. \pts{7}
\end{parts}

\medskip\rule{\linewidth}{0.2pt}

\medskip\noindent\textbf{Question 6.}
\begin{parts}
  \item Let the probability density function (p.d.f.) of a random variable $X$ be:
  $$f(x) = \begin{cases} ax, & 0 \le x \le 2 \\ a(3-x), & 2 < x < 3 \\ 0, & \text{elsewhere} \end{cases}$$
  Determine the value of $a$ and also compute $P(X < 1.5)$. \pts{7}
  \item The joint p.d.f. of a two-dimensional random variable $(X, Y)$ is given by:
  $$f(x, y) = \begin{cases} 8xy, & 0 < x < 1, \ 0 < y < x \\ 0, & \text{otherwise} \end{cases}$$
  Obtain the conditional p.d.f.s of $X$ given $Y=y$ and $Y$ given $X=x$. Also check for the independence of $X$ and $Y$. \pts{7}
\end{parts}

\medskip\rule{\linewidth}{0.2pt}

\medskip\noindent\textbf{Question 7.}
\begin{parts}
  \item Let the joint p.d.f. of $X$ and $Y$ be:
  $$f(x, y) = \begin{cases} e^{-(x+y)}, & x > 0, \ y > 0 \\ 0, & \text{elsewhere} \end{cases}$$
  Calculate $P(1 < X+Y < 2)$. \pts{7}
  \item If $X_1, X_2, \dots, X_n$ be any $n$ random variables and $a_1, a_2, \dots, a_n$ be any $n$ constants, then find the value of $\text{Var}\left(\sum_{i=1}^n a_i X_i\right)$. \pts{7}
\end{parts}

\medskip\rule{\linewidth}{0.2pt}

\medskip\noindent\textbf{Question 8.}
\begin{parts}
  \item State and prove the addition theorem of expectation for any $n$ jointly distributed random variables $Z_1, Z_2, \dots, Z_n$. \pts{7}
  \item The joint probability distribution of two random variables $X$ and $Y$ is given by:
  $$P(X=0, Y=-1) = \frac{1}{2}, \quad P(X=1, Y=-1) = \frac{1}{4}, \quad P(X=1, Y=1) = \frac{1}{4}$$
  Then find:
  \begin{parts}
    \item the marginal probability distributions of $X$ and $Y$;
    \item the conditional probability distribution of $X$ given $Y = 1$.
  \end{parts} \pts{7}
\end{parts}

\vfill
\rule{\linewidth}{0.4pt}
\begin{center}\small
  \href{https://bhu-exam.vercel.app/}{Click here to visit our website}\\[4pt]\href{https://me-aryan.vercel.app/}{Click here to visit our contributor}\\[4pt]\href{https://quashan-me.vercel.app/}{Click here to visit our contributor}
\end{center}

\end{document}